\documentclass[twoside,11pt]{article}

% JMLR style — exact file from the "Stochastic Interpolants" paper (arXiv:2303.08797).
% jmlr2e already loads amssymb, graphicx, natbib, hyperref — don't reload those.
\usepackage{jmlr2e}

\usepackage[T1]{fontenc}
\usepackage[utf8]{inputenc}
\usepackage[french,english]{babel}
\usepackage{amsmath}

% JMLR running headers
\ShortHeadings{Manuscrit}{Ripoll}
\firstpageno{1}

\title{Manuscrit}

\author{\name jules.ripoll \email jules.ripoll@insa-toulouse.fr \\
	\addr INSA Toulouse}
\editor{}

\begin{document}

\maketitle

\section{Flow-matching and Diffusion under the lense of Stochastic interpolants}

Tout ça c'est la même chose. Montrer qu'en partant du formalisme général des interpolants stochastiques on retombe à la fois sur les algos discrets de flow matching et ddpm etc...

\section{Le DiT en détails}

\section{Pourquoi est-ce que d'un coup c'est devenu ok de faire du image-space à l'échelle on a plus besoin de compression ?}

\section{Les tâches d'édition d'image sont mal définies}

\section{notes random}

Parler qlq part des méthodes de gatys et de spliceViT ? Ces papiers sont très cool

\bibliographystyle{plainnat}
\bibliography{references}

\end{document}
